% Topology graph with an attacker-controlled node and a defense boundary.
\documentclass[border=3pt]{standalone}
\usepackage[T1]{fontenc}
\usepackage{lmodern}
\usepackage{amsmath}
\usepackage{crypto-tikz}
\begin{document}
\begin{tikzpicture}[x=1cm,y=1cm]
  \node[benign, minimum width=13mm] (a) at (0,0) {A};
  \node[benign, minimum width=13mm] (b) at (2.2,1.2) {B};
  \node[defense, minimum width=13mm] (c) at (4.4,0) {C};
  \node[benign, minimum width=13mm] (d) at (2.2,-1.2) {D};
  \node[attacker, minimum width=16mm] (x) at (6.8,1.2) {external};
  \draw[wire] (a) -- (b) -- (c) -- (d) -- (a);
  \draw[msg] (x) -- node[above, font=\footnotesize] {injected message} (c);
  \draw[flow] (b) -- (c);
  \draw[flow] (d) -- (c);
  \node[op, anchor=west, align=left, fill=white, inner sep=2pt] at (-0.6,-2.35)
    {blue: peer handoff\\orange: attacker-controlled path\\outlined: defense boundary};
\end{tikzpicture}
\end{document}
